\subsection{从同步成本到时空结构：因果锥、传播上界与连续极限}\label{subsec:cdim-sync-causal-holonomy-spacetime}
上一小节已经把相位、射影态与干涉压回到局部--全局粘接残差在统计表示层上的投影。
若把同一类残差改在同步、共同支撑与闭路复合层上读取，则时间、空间、因果与曲率也不必被当作外加几何舞台，而可被理解为协议内部可比性条件的几何重组。
本小节只抽取这一重组的最小结构：同步所需的最小延迟、共同支撑所需的最小代价、以及闭路 residual 的面积归一化极限。
因此，这里既不预设连续流形，也不预设 Lorentz 度量或固定光速常数；所谓传播上界与连续曲率，只作为这些协议量之间的表示后果出现。

\subsubsection{同步延迟、共同支撑代价与因果前序}
设 \(\mathcal X\) 为协议允许的局部状态类。
对 \(x,y\in\mathcal X\)，记 \(x\leadsto y\) 为一条允许的单步精化、同步或传输步。

\begin{definition}[同步延迟与共同支撑代价]\label{def:cdim-sync-delay-support-cost}
对任意 \(x,y\in\mathcal X\)，定义：
\begin{enumerate}
  \item \textbf{同步延迟}
  \[
  \tau_{\mathrm{sync}}(x,y)\in[0,\infty]
  \]
  为把 \(x\) 与 \(y\) 拉入同一可比较协议口径所需的最小在线延迟、最小精化深度或最小同步成本；
  \item \textbf{共同支撑代价}
  \[
  d_{\mathrm{space}}(x,y)\in[0,\infty]
  \]
  为把 \(x\) 与 \(y\) 拉入共同支撑、共同 forcing 区域或共同可比较局部图所需的最小资源代价。
\end{enumerate}
若不存在任何允许程序使该比较得以建立，则对应量记为 \(+\infty\)。
\end{definition}

\begin{definition}[因果可达前序]\label{def:cdim-causal-reachability-preorder}
对 \(x,y\in\mathcal X\)，若存在有限允许链
\[
x=x_0\leadsto x_1\leadsto \cdots \leadsto x_n=y
\]
使得
\[
\sum_{k=0}^{n-1}\tau_{\mathrm{sync}}(x_k,x_{k+1})<\infty,
\qquad
\sum_{k=0}^{n-1}d_{\mathrm{space}}(x_k,x_{k+1})<\infty,
\]
则记
\[
x\preceq y,
\]
称 \(y\) 落在 \(x\) 的\textbf{因果可达域}中。
\end{definition}

\begin{proposition}[同步与共同支撑诱导的因果前序]\label{prop:cdim-causal-preorder-from-subadditivity}
若同步延迟与共同支撑代价沿允许复合满足
\[
\tau_{\mathrm{sync}}(x,z)
\le
\tau_{\mathrm{sync}}(x,y)+\tau_{\mathrm{sync}}(y,z),
\]
\[
d_{\mathrm{space}}(x,z)
\le
d_{\mathrm{space}}(x,y)+d_{\mathrm{space}}(y,z),
\]
且在对角线上均为 \(0\)，则定义 \ref{def:cdim-causal-reachability-preorder} 所给出的关系 \(\preceq\) 构成 \(\mathcal X\) 上的一个前序。
\leanverified{paper\_cdim\_causal\_preorder\_from\_subadditivity}
\end{proposition}
\begin{proof}
反身性由零长度链给出。
若 \(x\preceq y\) 且 \(y\preceq z\)，则分别存在从 \(x\) 到 \(y\) 与从 \(y\) 到 \(z\) 的有限允许链。
将两条链首尾拼接，再利用同步延迟与共同支撑代价的次可加性，即得一条从 \(x\) 到 \(z\) 的有限允许链，其总同步成本与总支撑成本仍为有限，故 \(x\preceq z\)。
证毕。
\end{proof}

\begin{remark}[当前稿件中的同步接口]\label{rem:cdim-sync-delay-existing-interfaces}
定义 \ref{def:cdim-sync-delay-support-cost} 并非空泛记号，而已有若干后端实现。
在折叠同步模型中，最坏同步延迟由推论 \ref{cor:Ym-sync-delay} 精确给出为 \(D_{\mathrm{sync}}(m)=m-1\)；在因果右逆口径下，最小延迟与最小状态数由定理 \ref{thm:xi-time-min-delay-min-states} 同时取到；在固定窗口的概率同步口径下，最小精确同步成本又可由定理 \ref{thm:conclusion-window6-exact-fold-synchronization-mutual-information} 写成互信息极小值。
因此，同步量在本文中已经具备组合论、自动机与信息论三种可核验落点。
\end{remark}

\subsubsection{传播上界作为同步--支撑比的极大不变量}
若一条允许步既伴随有限同步延迟，又伴随有限共同支撑代价，则其空间--时间比给出该步的局部传播速率。

\begin{definition}[局部传播比与传播上界]\label{def:cdim-local-propagation-ratio}
对任一允许的基本传播步
\[
e:x\leadsto y
\qquad
\text{且}\qquad
0<\tau_{\mathrm{sync}}(x,y)<\infty,
\]
定义其\textbf{局部传播比}为
\[
\rho(e)
:=
\frac{d_{\mathrm{space}}(x,y)}{\tau_{\mathrm{sync}}(x,y)}.
\]
进一步定义
\[
c_*:=\sup\bigl\{\rho(e): e\ \text{为允许的基本传播步}\bigr\}\in[0,\infty].
\]
称 \(c_*\) 为该协议族的\textbf{传播上界}。
\end{definition}

\begin{proposition}[链式传播上界]\label{prop:cdim-chainwise-propagation-bound}
设
\[
x=x_0\leadsto x_1\leadsto\cdots\leadsto x_n=y
\]
为任一允许链，并定义其总同步延迟与总共同支撑代价
\[
\tau_{\mathrm{chain}}
:=
\sum_{k=0}^{n-1}\tau_{\mathrm{sync}}(x_k,x_{k+1}),
\qquad
d_{\mathrm{chain}}
:=
\sum_{k=0}^{n-1}d_{\mathrm{space}}(x_k,x_{k+1}).
\]
若定义 \ref{def:cdim-local-propagation-ratio} 中的 \(c_*\) 有限，且链上每一步满足
\[
\tau_{\mathrm{sync}}(x_k,x_{k+1})>0
\quad\text{或}\quad
d_{\mathrm{space}}(x_k,x_{k+1})=0,
\]
则恒有
\[
d_{\mathrm{chain}}\le c_*\,\tau_{\mathrm{chain}}.
\]
\leanverified{paper\_cdim\_chainwise\_propagation\_bound}
\end{proposition}
\begin{proof}
由 \(c_*\) 的定义，对每一步 \(x_k\leadsto x_{k+1}\) 都有
\[
d_{\mathrm{space}}(x_k,x_{k+1})
\le
c_*\,\tau_{\mathrm{sync}}(x_k,x_{k+1}).
\]
对 \(k\) 求和即得所示不等式。
\end{proof}

\begin{remark}[局域可构造边与全局锥边界的分层]\label{rem:cdim-local-edge-vs-global-cone-layering}
这里需要区分两种不同层级的几何对象。单条允许步 \(x\leadsto y\) 及其同步延迟、共同支撑代价，给出的是局域可构造边；它只说明协议内部存在一条合法传输或同步步骤。相反，\(c_*\)、\(J^\pm\)、\(I^\pm\) 与 \(N^\pm\) 则属于全局刚性层，因为它们要在所有允许步与所有有限链上取上确界、闭包与边界。也就是说，因果锥边界不是一根先验画好的线，而是从大量局域可构造边经比例控制与链式闭包后才显现出来的全局不变量。
\end{remark}

\begin{definition}[未来锥、类时区与边界候选]\label{def:cdim-causal-cone-candidates}
对任意 \(x\in\mathcal X\)，定义其未来与过去锥候选为
\[
J^+(x):=\{y\in\mathcal X:\ x\preceq y\},
\qquad
J^-(x):=\{y\in\mathcal X:\ y\preceq x\}.
\]
若存在一条从 \(x\) 到 \(y\) 的允许链以及某个 \(\varepsilon>0\)，使
\[
d_{\mathrm{chain}}
\le
(c_*-\varepsilon)\tau_{\mathrm{chain}},
\]
则称 \(y\) 属于 \(x\) 的\textbf{类时未来候选}，记为 \(I^+(x)\)；并定义传播边界候选
\[
N^+(x):=J^+(x)\setminus I^+(x).
\]
过去类时区与边界 \(I^-(x),N^-(x)\) 同理定义。
\end{definition}

\begin{conjecture}[传播锥的观察者协变性]\label{conj:cdim-causal-cone-observer-covariance}
设 \(T:\mathcal X\to\mathcal X\) 为一个允许的观察者变换，并满足：
\begin{enumerate}
  \item \(T\) 把允许的局部传播步送到允许的局部传播步；
  \item \(T\) 保持同步延迟与共同支撑代价在同一规范账本下的等价类；
  \item \(T\) 保持局部细化、粗粒化与张量复合的可比较性结构。
\end{enumerate}
则预期有：
\begin{enumerate}
  \item \(x\preceq y\) 当且仅当 \(Tx\preceq Ty\)；
  \item 传播上界 \(c_*\) 在 \(T\) 下不变；
  \item \(T(J^\pm(x))=J^\pm(Tx)\)、\(T(I^\pm(x))=I^\pm(Tx)\)、\(T(N^\pm(x))=N^\pm(Tx)\)。
\end{enumerate}
换言之，观察者变化首先应保持的不是某个既定坐标表达，而是因果可达性、锥内部与锥边界的分层。
\end{conjecture}

\subsubsection{小回路 holonomy、归一化离散曲率与连续极限}
若上一小节中的可见相位来自局部 lift 在重叠上的粘接残差，则几何侧的曲率应当来自同一类残差沿小回路复合后的累积。
在本文当前稿件中，这一点已经在 dyadic holonomy 的边界重建与通量缺陷理论里获得了严格后端。

\begin{definition}[dyadic 归一化 holonomy 曲率场]\label{def:cdim-dyadic-normalized-holonomy-curvature}
沿用定理 \ref{thm:conclusion-nonabelian-dyadic-stokes-energy-reconstruction} 的设定。
对边长 \(h_m=2^{-m}\) 的 dyadic 方形剖分 \(\mathcal Q_m\)，若
\[
\operatorname{Hol}_A(\partial Q)=\exp(\Omega(Q))
\qquad (Q\in\mathcal Q_m),
\]
则定义分片常值的\textbf{归一化离散曲率场}
\[
\mathscr R_m(x)
:=
\sum_{Q\in\mathcal Q_m}
\frac{\Omega(Q)}{h_m^2}\,\mathbf 1_Q(x).
\]
它把小回路 holonomy 的对数主项写成面积归一化后的局部曲率密度。
\end{definition}

\begin{corollary}[dyadic holonomy 归一化曲率的连续极限]\label{cor:cdim-dyadic-holonomy-curvature-limit}
在定理 \ref{thm:conclusion-nonabelian-dyadic-stokes-energy-reconstruction} 的设定下，有
\[
\mathscr R_m\longrightarrow F_{12}
\qquad\text{于 }L^2(R_1),
\]
并且
\[
\|\mathscr R_m\|_{L^2(R_1)}^2
\longrightarrow
\int_{R_1}|F_{12}(x)|^2\,\mathrm d x.
\]
因此，连续曲率对象在该 dyadic 后端上可被解释为归一化小回路 holonomy 的 \(L^2\)-极限。
\end{corollary}
\begin{proof}
在定理 \ref{thm:conclusion-nonabelian-dyadic-stokes-energy-reconstruction} 的证明中，\(\mathscr R_m\) 正是其中记作 \(\mathscr H_m\) 的分片常值场。
该证明已经给出
\[
\|\mathscr R_m-P_mF_{12}\|_{L^2(R_1)}\longrightarrow 0,
\]
其中 \(P_mF_{12}\) 为 \(F_{12}\) 关于 dyadic \(\sigma\)-代数的条件期望。
又 \(P_mF_{12}\to F_{12}\) 于 \(L^2(R_1)\)，故由三角不等式得
\[
\mathscr R_m\to F_{12}\qquad\text{于 }L^2(R_1).
\]
第二式即定理 \ref{thm:conclusion-nonabelian-dyadic-stokes-energy-reconstruction} 对 \(\mathcal E_m^{\mathrm{NA}}(A)=\|\mathscr R_m\|_2^2\) 的陈述。
\end{proof}

\leanverified{paper\_cdim\_dyadic\_holonomy\_curvature\_limit}

\begin{remark}[通量缺陷给出的后验证书]\label{rem:cdim-dyadic-holonomy-flux-defect-certificate}
定理 \ref{thm:conclusion-nonabelian-dyadic-flux-defect-telescoping} 进一步表明：对任意 dyadic 方形 \(Q_0\)，粗尺度 holonomy 对数 \(\Omega(Q_0)\) 与连续曲率通量 \(\int_{Q_0}F_{12}\) 之差，可以完全重写为一列绝对收敛的 dyadic 通量缺陷级数。
因此，连续曲率不仅是小回路 holonomy 的极限对象，而且其偏离量本身也可由纯边界数据驱动的多尺度证书链后验审计。
\end{remark}

\begin{conjecture}[holonomy 细化极限的普适曲率候选]\label{conj:cdim-holonomy-curvature-refinement-limit}
设 \(\{\mathcal U_h\}_{h\downarrow 0}\) 为一族细化覆盖，且每个尺度 \(h\) 上都给定局部过渡数据与闭路 holonomy，使得：
\begin{enumerate}
  \item 小回路对数 holonomy 在尺度 \(h\) 上具有一致的主支选择与局部线性化；
  \item 面积归一化后的离散曲率场在紧集上弱预紧；
  \item 离散 Bianchi 型约束与细化过程相容；
  \item 多尺度通量缺陷满足与定理 \ref{thm:conclusion-nonabelian-dyadic-flux-defect-telescoping} 同型的可求和控制。
\end{enumerate}
则应存在子列 \(h_j\downarrow 0\)，使相应的归一化离散曲率场弱收敛到某个连续曲率候选 \(R\)；并且若极限唯一，则 \(R\) 只由离散 holonomy 的细化极限决定。
这表明连续曲率不必预先给定，而应当作为局部粘接残差在小回路复合层上的尺度稳定首项出现。
\end{conjecture}

\noindent
因此，在本章当前语境下，时间、空间、因果与曲率都可以被压回到同一条协议链上：同步延迟决定时间尺度，共同支撑代价决定空间比较，二者的极大传播比决定因果锥的边界，而小回路 holonomy 的面积归一化极限则给出连续曲率候选。
这条链仍未恢复唯一的 Lorentz 型度量，也未导出场方程；但它已经说明，所谓时空骨架并不位于协议之外，而是由同步、传播、比较与闭路 residual 共同选出的几何组织。
在本章的后续编译接口中，这也意味着地址对象不能只记录单一可见前缀，而必须继续把可见轴、残差轴与模式轴一起显式化，以承载同步账本、传播约束与 residual 的可核验外置。

\endinput
